\documentclass[12pt]{article}
\usepackage[utf8]{inputenc}
\usepackage[T1]{fontenc}
\usepackage{hyperref}
\title{Tsimihety Malagasy Language Resources for K-12 Curriculum Development}
\author{Compiled by Claude Sonnet 5.5 (claude-sonnet-5-5) for Andr\'as Kornai}
\date{}
\begin{document}
\maketitle

\section{Introduction}
Tsimihety Malagasy (ISO 639-3 xmw; Glottocode tsim1257) is a northern Malagasy variety spoken mainly in the Sofia Region of Madagascar, with diaspora communities in Mahajanga and Antananarivo. My searches found no Tsimihety-specific school materials. The resources below are mostly reference and lexical.

\section{Primary Educational Resources}
No Tsimihety primers, textbooks or school programmes were found in search results.

\section{Linguistic and Reference Materials}
\begin{itemize}
\item Wikipedia, ``Tsimihety dialect'', with a vocabulary comparison against Standard Malagasy: \url{https://en.wikipedia.org/wiki/Tsimihety_dialect}
\item Adelaar, ``Malagasy Dialect Divisions: Genetic versus Emblematic Criteria'' (University of Melbourne repository; I identified the title from the search snippet only): \url{https://minerva-access.unimelb.edu.au/server/api/core/bitstreams/15b6ed76-ba89-50d2-ab78-366ce14ddc35/content}
\item ``Malagasy Dialects and the Peopling of Madagascar'', whose consultants include a Tsimihety speaker from Mandritsara: \url{https://arxiv.org/pdf/1102.2180}
\end{itemize}

\section{Religious and Community Resources}
No Tsimihety-specific scripture or community organisation was found.

\section{Digital and Audio Resources}
\begin{itemize}
\item Malagasyword.org, Tsimihety entry with dialectal words: \url{https://malagasyword.org/bins/teny2/Tsimihety}
\item Malagasyword.org entry tagged as Tsimihety usage: \url{https://malagasyword.org/bins/teny2/mpis\^ova}
\end{itemize}
The site's bibliography is said to list a Tsimihety dictionary, \emph{Rantimbola\~na}, but I did not confirm that it is available online.

\section{Curriculum Development Context}
Tsimihety has a large speaker base, but I found no teaching materials for it, and it is not known to be taught in schools. A curriculum would need to build on Standard Malagasy materials and on community-based documentation of Tsimihety.
\end{document}